% Part: normal-modal-logic
% Chapter: axioms-systems
% Section: normal-logics

\documentclass[../../../include/open-logic-section]{subfiles}

\begin{document}

\olfileid{nml}{prf}{nor}

\olsection{正规模态逻辑}

并非每个模态公式集都能轻易地刻画为从一组公理可推出的公式集。我们希望模态逻辑具有良好的性质。首先，我们在经典命题逻辑中能推出的一切应该仍然可推出，当然，此时要考虑到公式现在可能还包含 \iftag{prvBox}{$\Box$\iftag{prvDiamond}{ 和~}{}}{}\iftag{prvDiamond}{$\Diamond$}{}。为此，我们要求一个模态逻辑包含所有重言式实例，并且对肯定前件封闭。

\begin{defn}
  一个\emph{模态逻辑}是满足以下条件的模态公式集合~$\Sigma$：
  \begin{enumerate}
  \item 包含所有重言式，并且
  \item 对代入封闭，即如果 $!A \in \Sigma$，且 $!D_1$, \dots, $!D_n$ 是公式，则
    \[
    \SSubst{!A}{\subst{!D_1}{p_1}, \dots, \subst{!D_n}{p_n}} \in \Sigma,
    \]
    \item 对\emph{肯定前件}封闭，即如果 $!A \in \Sigma$ 且 $!A \lif !B \in \Sigma$，则 $!B \in \Sigma$。
  \end{enumerate}
\end{defn}

为了将关系语义用于模态逻辑，我们还必须要求它包含所有在所有模态模型中都有效的公式。事实证明，只要 \Ax{K}\iftag{prvDiamond}{ 和~\Dual{}}{} 的所有实例都是可推出的，并且每当公式~$!A$ 可推出时~$\Box !A$ 也可推出，这一要求即可满足。满足这些条件的模态逻辑称为\emph{正规的}。（当然，也存在非正规的模态逻辑，但通常的关系模型并不适用于它们。）

\begin{defn}
  如果模态逻辑 $\Sigma$ 包含
  \begin{align*}
    \tag{\Ax{K}} & \Box(p \lif q) \lif (\Box p \lif \Box q),
    \iftag{prvDiamond}{\\
      \tag{\Dual} & \Diamond p \liff \lnot\Box\lnot p}{}
  \end{align*}
  并且对\emph{必然化}封闭，即如果 $!A \in \Sigma$，则 $\Box !A \in \Sigma$，则称它是\emph{正规的}。
\end{defn}

注意，虽然重言后承“足够精细”，能够保持\emph{世界中的真}，但规则 \Nec{} 仅保持\emph{模型中的真}（因而也保持框架或框架类中的有效性）。

\begin{prop}\ollabel{prop:rk}
  每个正规模态逻辑都对规则 \RK 封闭：
  \begin{prooftree}
    \AxiomC{$!A_1 \lif (!A_2 \lif \cdots (!A_{n-1} \lif !A_n)\cdots)$}
    \RightLabel{\RK}
    \UnaryInfC{$\Box!A_1 \lif (\Box!A_2 \lif \cdots (\Box!A_{n-1}
      \lif \Box!A_n)\cdots).$}
  \end{prooftree}
\end{prop}

\begin{proof}
  对~$n$ 进行归纳：若 $n = 1$，则该规则正是 \Nec，而每个正规模态逻辑都对 \Nec 封闭。

  现在假设结论对 $n-1$ 成立；我们证明它对~$n$ 成立。

  假设
  \begin{align*}
  & !A_1 \lif (!A_2 \lif \cdots (!A_{n-1} \lif !A_n)\cdots) \in \Sigma
  \intertext{由归纳假设，我们有}
  & \Box!A_1 \lif (\Box!A_2 \lif \cdots \Box(!A_{n-1} \lif !A_n)\cdots)
  \in \Sigma
  \intertext{由于 $\Sigma$ 是正规模态逻辑，它包含~$\Ax{K}$ 的所有实例，特别地有}
  & \Box(!A_{n-1} \lif !A_n) \lif (\Box!A_{n-1} \lif \Box!A_n) \in \Sigma
  \intertext{利用肯定前件和适当的重言式实例，我们得到}
  & \Box!A_1 \lif (\Box!A_2 \lif \cdots (\Box!A_{n-1}
  \lif \Box!A_n)\cdots) \in \Sigma.
  \end{align*}
\end{proof}

\begin{prop}\ollabel{prop:notDiamondBot}
  每个正规模态逻辑 $\Sigma$ 都包含~$\lnot\Diamond\lfalse$。
\end{prop}

\begin{prob}
  证明\olref[nml][prf][nor]{prop:notDiamondBot}。
\end{prob}

\begin{prop}
  设 $!A_1$, \dots, $!A_n$ 为公式。则存在包含 $!A_1$, \dots, $!A_n$ 的所有实例的最小模态逻辑 $\Sigma$。
\end{prop}

\begin{proof}
  给定 $!A_1$, \dots, $!A_n$，将 $\Sigma$ 定义为所有包含 $!A_1$, \dots, $!A_n$ 的所有实例的正规模态逻辑的交集。该交集非空，因为所有公式的集合 $\Frm[L]$ 本身就是这样一个模态逻辑。
\end{proof}

\begin{defn}
包含 $!A_1$, \dots, $!A_n$ 的最小正规模态逻辑称为一个\emph{模态系统}，记作 $\Log{K} !A_1 \dots
!A_n$。最小的正规模态逻辑记作~\Log{K}。
\end{defn}

\end{document}